\begin{longtable}{|C{1.4cm}|L{4.0cm}|L{9.4cm}|}
\hline
\thd{RN} & \thd{NOMBRE DE LA REGLA} & \thd{DETALLE DE LA REGLA} \\
\hline
RN-001 & Conjunto cerrado de roles & \textbf{SI} una cuenta opera en la plataforma, \textbf{ENTONCES} el sistema le reconoce exactamente uno de dos roles ---SuperAdmin o Administrador--- y aplica únicamente los permisos de ese rol. \textbf{SE APLICA A} todas las cuentas y a todos los módulos y endpoints de la plataforma. \textbf{SE BASA EN} el modelo RBAC de dos roles acordado con el cliente (acta CS3081-004-2026 · REQ-10) y en el principio de mínimo privilegio. \textbf{Y TIENE LA SIGUIENTE EXCEPCIÓN} no se admiten roles adicionales (p. ej. «Usuario») ni cuentas sin rol asignado. \\
\hline
RN-002 & Unicidad del SuperAdmin & \textbf{SI} se crea, habilita, deshabilita o transfiere una cuenta, \textbf{ENTONCES} el sistema garantiza que exista exactamente una cuenta con rol SuperAdmin en todo momento, con independencia de su estado de habilitación o de si mantiene una sesión abierta. \textbf{SE APLICA A} la creación de cuentas, la habilitación/deshabilitación y la transferencia de rol. \textbf{SE BASA EN} el principio de responsabilidad única sobre la administración del sistema y en el acuerdo del cliente (acta CS3081-004-2026 · REQ-10). \textbf{Y TIENE LA SIGUIENTE EXCEPCIÓN} la cuenta SuperAdmin inicial se crea durante el despliegue y no puede crearse desde la aplicación (RN-006). \\
\hline
RN-003 & Permisos cerrados por rol & \textbf{SI} una cuenta intenta ejecutar una acción, \textbf{ENTONCES} el sistema la autoriza únicamente si esa acción corresponde a su rol; en caso contrario, deniega el acceso con error de autorización. \textbf{SE APLICA A} todos los módulos y endpoints, evaluándose en cada petición. \textbf{SE BASA EN} el principio de mínimo privilegio (RBAC). \textbf{Y TIENE LA SIGUIENTE EXCEPCIÓN} no existe excepción por interfaz: ocultar una opción en la pantalla no otorga ni niega permiso, pues la autorización se resuelve en el servidor. \\
\hline
RN-004 & Condiciones de autenticación & \textbf{SI} una cuenta intenta iniciar sesión, \textbf{ENTONCES} el sistema le permite el acceso únicamente si está registrada, se encuentra habilitada y su contraseña es correcta. \textbf{SE APLICA A} todas las cuentas, incluida la cuenta SuperAdmin creada en el despliegue. \textbf{SE BASA EN} el control de acceso por credenciales. \textbf{Y TIENE LA SIGUIENTE EXCEPCIÓN} una cuenta deshabilitada no puede autenticarse aunque sus credenciales sean correctas (RN-005). \\
\hline
RN-005 & Efecto de la deshabilitación de cuenta & \textbf{SI} el SuperAdmin deshabilita una cuenta, \textbf{ENTONCES} el sistema invalida de inmediato todas sus sesiones activas y le revoca el acceso a todas las funcionalidades, sin eliminar su información ni su historial. \textbf{SE APLICA A} las cuentas de Administrador. \textbf{SE BASA EN} la revocación inmediata de accesos y en la conservación de la trazabilidad. \textbf{Y TIENE LA SIGUIENTE EXCEPCIÓN} la cuenta SuperAdmin no puede deshabilitarse desde la aplicación; primero debe transferirse el rol (RN-009). \\
\hline
RN-006 & Cuenta SuperAdmin inicial & \textbf{SI} el sistema se despliega por primera vez, \textbf{ENTONCES} se crea una única cuenta con rol SuperAdmin durante el despliegue, con independencia de la aplicación. \textbf{SE APLICA A} el aprovisionamiento inicial del sistema. \textbf{SE BASA EN} la necesidad de un administrador raíz que no dependa de la interfaz. \textbf{Y TIENE LA SIGUIENTE EXCEPCIÓN} esta cuenta no puede crearse ni eliminarse desde la aplicación; su rol solo se transfiere (RN-009). \\
\hline
RN-007 & Datos obligatorios de la cuenta & \textbf{SI} se registra o edita una cuenta, \textbf{ENTONCES} el sistema exige nombre completo, correo electrónico y contraseña, y valida el formato del correo. \textbf{SE APLICA A} la creación y la edición de cuentas. \textbf{SE BASA EN} la integridad de la identidad de la cuenta. \textbf{Y TIENE LA SIGUIENTE EXCEPCIÓN} no se admite una cuenta sin los tres datos obligatorios ni con un correo ya registrado (unicidad de correo). \\
\hline
RN-008 & Rol por defecto & \textbf{SI} el SuperAdmin crea una cuenta desde la aplicación, \textbf{ENTONCES} el sistema le asigna automáticamente el rol Administrador. \textbf{SE APLICA A} la creación de cuentas. \textbf{SE BASA EN} que el rol SuperAdmin solo se obtiene por transferencia (RN-009). \textbf{Y TIENE LA SIGUIENTE EXCEPCIÓN} no es posible crear cuentas con rol SuperAdmin desde la interfaz. \\
\hline
RN-009 & Transferencia atómica del rol SuperAdmin & \textbf{SI} el SuperAdmin transfiere su rol a una cuenta de Administrador existente y habilitada, \textbf{ENTONCES} el sistema ejecuta la operación de forma atómica, previa confirmación: la cuenta destino adquiere el rol SuperAdmin y la cuenta origen pasa a Administrador en una sola transacción. \textbf{SE APLICA A} la transferencia de rol. \textbf{SE BASA EN} la unicidad del SuperAdmin (RN-002) y en la integridad transaccional. \textbf{Y TIENE LA SIGUIENTE EXCEPCIÓN} si la cuenta destino no existe, está deshabilitada o ya es SuperAdmin, la transferencia se rechaza sin alterar el estado de las cuentas. \\
\hline
RN-010 & Autonomía en la recuperación de acceso & \textbf{SI} un usuario pierde sus credenciales, \textbf{ENTONCES} el sistema le permite recuperar el acceso por sí mismo mediante un enlace seguro enviado a su correo, sin intervención del SuperAdmin. \textbf{SE APLICA A} todas las cuentas habilitadas. \textbf{SE BASA EN} la autonomía del usuario y en la autenticación segura por correo. \textbf{Y TIENE LA SIGUIENTE EXCEPCIÓN} el enlace es de un solo uso y de vigencia limitada (30 minutos); vencido o ya usado, debe generarse uno nuevo. \\
\hline
RN-011 & Política de contraseñas & \textbf{SI} se crea o cambia una contraseña, \textbf{ENTONCES} el sistema exige un mínimo de 8 caracteres con al menos una mayúscula, una minúscula, un dígito y un carácter especial, y no admite que coincida con el correo de la cuenta. \textbf{SE APLICA A} el registro, el cambio y el restablecimiento de contraseña. \textbf{SE BASA EN} las buenas prácticas de autenticación segura. \textbf{Y TIENE LA SIGUIENTE EXCEPCIÓN} no se permite guardar la contraseña en texto plano; su almacenamiento es un hash con sal (RNF). \\
\hline
RN-012 & Bloqueo por intentos fallidos & \textbf{SI} una cuenta acumula 5 intentos fallidos consecutivos de inicio de sesión, \textbf{ENTONCES} el sistema la bloquea temporalmente durante 15 minutos, impide nuevos intentos y registra el evento en la auditoría. \textbf{SE APLICA A} todas las cuentas que se autentican con credenciales. \textbf{SE BASA EN} la mitigación de ataques de fuerza bruta. \textbf{Y TIENE LA SIGUIENTE EXCEPCIÓN} el bloqueo se levanta automáticamente al vencer el tiempo o al restablecerse la contraseña; no requiere intervención manual. \\
\hline
RN-013 & Sesión e inactividad & \textbf{SI} una sesión permanece sin actividad durante 2 horas, \textbf{ENTONCES} el sistema la expira y exige una nueva autenticación; asimismo, el usuario puede cerrar la sesión manualmente. \textbf{SE APLICA A} todas las sesiones autenticadas. \textbf{SE BASA EN} el control de sesiones y en la revocación por inactividad. \textbf{Y TIENE LA SIGUIENTE EXCEPCIÓN} tras una transferencia de rol, la sesión activa aplica el rol vigente; si la cuenta se deshabilita, la sesión se invalida de inmediato (RN-005). \\
\hline
RN-014 & Registro de empresas & \textbf{SI} el SuperAdmin registra una empresa, \textbf{ENTONCES} el sistema exige un nombre y la asignación a un sector. \textbf{SE APLICA A} la creación y edición del catálogo de empresas. \textbf{SE BASA EN} el acuerdo con el cliente (acta CS3081-004-2026 · REQ-10) y en la necesidad de identificar la empresa por su nombre y su sector. \textbf{Y TIENE LA SIGUIENTE EXCEPCIÓN} no se admite una empresa sin nombre ni sin sector. \\
\hline
RN-015 & Sectores habilitados (conjunto cerrado) & \textbf{SI} se registra una empresa o se realiza una ingesta o consulta, \textbf{ENTONCES} el sistema solo admite tres sectores: Minería, Petróleo y Gas, y Energía. \textbf{SE APLICA A} el catálogo de empresas, la ingesta y las consultas. \textbf{SE BASA EN} el acuerdo con el cliente (acta CS3081-003-2026 · REQ-07). \textbf{Y TIENE LA SIGUIENTE EXCEPCIÓN} no se crean sectores nuevos; las empresas de otros sectores pueden registrarse, pero no son elegibles para el análisis ni para la generación de reportes. \\
\hline
RN-016 & Unicidad de la empresa & \textbf{SI} se registra una empresa, \textbf{ENTONCES} el sistema rechaza el registro si ya existe otra con el mismo nombre. \textbf{SE APLICA A} la creación y edición del catálogo de empresas. \textbf{SE BASA EN} la integridad del catálogo. \textbf{Y TIENE LA SIGUIENTE EXCEPCIÓN} no se admiten empresas duplicadas. \\
\hline
RN-017 & Vigencia de la empresa (activa/inactiva) & \textbf{SI} una empresa se marca como inactiva, \textbf{ENTONCES} el sistema deja de ofrecerla para la ingesta y las consultas y no la considera en el análisis. \textbf{SE APLICA A} el catálogo de empresas. \textbf{SE BASA EN} la vigencia del catálogo. \textbf{Y TIENE LA SIGUIENTE EXCEPCIÓN} la empresa inactiva conserva sus documentos y su historial y no se elimina. \\
\hline
RN-018 & Formato de ingesta & \textbf{SI} el SuperAdmin carga un documento, \textbf{ENTONCES} el sistema solo admite archivos Markdown (.md) con tablas de pipes. \textbf{SE APLICA A} la ingesta de memorias anuales y reportes de sostenibilidad. \textbf{SE BASA EN} el acuerdo del cliente (acta CS3081-005-2026 · REQ-20). \textbf{Y TIENE LA SIGUIENTE EXCEPCIÓN} el sistema no convierte formatos distintos de Markdown. \\
\hline
RN-019 & Responsabilidad de la conversión & \textbf{SI} el cliente entrega un documento, \textbf{ENTONCES} es responsable de su conversión a Markdown y de la integridad de sus tablas (formato pipe). \textbf{SE APLICA A} la preparación de los documentos fuente. \textbf{SE BASA EN} la delimitación de responsabilidades acordada. \textbf{Y TIENE LA SIGUIENTE EXCEPCIÓN} el sistema no reconvierte ni corrige el contenido entregado. \\
\hline
RN-020 & Tipos de documento admisibles & \textbf{SI} se carga un documento, \textbf{ENTONCES} su tipo debe ser memoria anual o reporte de sostenibilidad GRI. \textbf{SE APLICA A} la ingesta. \textbf{SE BASA EN} la distinción de los dos documentos fuente del proceso. \textbf{Y TIENE LA SIGUIENTE EXCEPCIÓN} no se admiten otros tipos de documento. \\
\hline
RN-021 & Tamaño máximo de ingesta & \textbf{SI} se carga un documento, \textbf{ENTONCES} el sistema rechaza la carga si el archivo supera los 50 MB. \textbf{SE APLICA A} la ingesta. \textbf{SE BASA EN} el límite acordado para el procesamiento. \textbf{Y TIENE LA SIGUIENTE EXCEPCIÓN} el rechazo se registra con su motivo y no se procesa contenido parcial. \\
\hline
RN-022 & Identificación obligatoria del documento & \textbf{SI} se carga un documento, \textbf{ENTONCES} debe estar asociado a una empresa activa, al año del ejercicio y a su tipo. \textbf{SE APLICA A} la ingesta. \textbf{SE BASA EN} la trazabilidad del documento de origen. \textbf{Y TIENE LA SIGUIENTE EXCEPCIÓN} no se indexa un documento sin empresa, año y tipo. \\
\hline
RN-023 & Unicidad de documento por empresa, año y tipo & \textbf{SI} se carga un documento, \textbf{ENTONCES} el sistema rechaza la carga si ya existe otro del mismo tipo para la misma empresa y año. \textbf{SE APLICA A} la ingesta. \textbf{SE BASA EN} la evitación de duplicados documentales. \textbf{Y TIENE LA SIGUIENTE EXCEPCIÓN} el rechazo indica la fecha y la cuenta que realizó la carga original. \\
\hline
RN-024 & Unicidad por contenido (SHA-256) & \textbf{SI} se carga un documento, \textbf{ENTONCES} el sistema calcula su hash SHA-256 y rechaza la carga si ese contenido ya fue ingestado. \textbf{SE APLICA A} la ingesta. \textbf{SE BASA EN} la detección de duplicados por contenido. \textbf{Y TIENE LA SIGUIENTE EXCEPCIÓN} un documento rechazado o interrumpido no deja contenido en el corpus. \\
\hline
RN-025 & Contenido mínimo y estado observado & \textbf{SI} un documento no contiene al menos un código del catálogo GRI ni una mención de sanción reconocible, \textbf{ENTONCES} el sistema lo marca como observado y no lo indexa hasta corregirse. \textbf{SE APLICA A} la ingesta. \textbf{SE BASA EN} la necesidad de que el documento aporte información analizable. \textbf{Y TIENE LA SIGUIENTE EXCEPCIÓN} el documento observado no se rechaza por completo; se conserva para su reintento. \\
\hline
RN-026 & Integridad atómica de la ingesta & \textbf{SI} una ingesta es rechazada o interrumpida, \textbf{ENTONCES} el sistema revierte el procesamiento y elimina cualquier contenido, fragmento, embedding o resultado generado parcialmente, conservando únicamente el motivo del rechazo. \textbf{SE APLICA A} la ingesta. \textbf{SE BASA EN} la integridad del corpus. \textbf{Y TIENE LA SIGUIENTE EXCEPCIÓN} solo se conserva la información necesaria para identificar y auditar el rechazo. \\
\hline
RN-027 & Catálogo de códigos GRI evaluables & \textbf{SI} el sistema analiza un documento, \textbf{ENTONCES} solo considera los 40 códigos del catálogo GRI corporativo versionado. \textbf{SE APLICA A} el análisis de brechas. \textbf{SE BASA EN} el catálogo acordado con el PO. \textbf{Y TIENE LA SIGUIENTE EXCEPCIÓN} el texto del estándar GRI oficial no es fuente de evaluación; la evidencia válida es la reportada por la empresa. \\
\hline
RN-028 & Detección de códigos y extracción de cita & \textbf{SI} el sistema analiza un documento indexado, \textbf{ENTONCES} identifica los códigos del catálogo presentes y extrae la cita textual de respaldo, sin evaluar su contenido. \textbf{SE APLICA A} el análisis posterior a la ingesta. \textbf{SE BASA EN} la supervisión humana del estado. \textbf{Y TIENE LA SIGUIENTE EXCEPCIÓN} el sistema no compara el contenido contra el texto del estándar GRI. \\
\hline
RN-029 & Asignación manual del estado GRI & \textbf{SI} se detecta un código GRI en el documento, \textbf{ENTONCES} el estado del código es asignado manualmente por el Administrador tras revisar la cita textual. \textbf{SE APLICA A} el análisis y la validación de brechas. \textbf{SE BASA EN} la supervisión humana acordada con el PO. \textbf{Y TIENE LA SIGUIENTE EXCEPCIÓN} el sistema no infiere ni sugiere el estado de forma automática. \\
\hline
RN-030 & Conjunto cerrado de estados GRI & \textbf{SI} el Administrador asigna un estado, \textbf{ENTONCES} solo puede ser uno de tres valores: OK, Baja sustancia o Sub-reportado. \textbf{SE APLICA A} el análisis y la validación de brechas. \textbf{SE BASA EN} los estados acordados con el PO. \textbf{Y TIENE LA SIGUIENTE EXCEPCIÓN} no existe un cuarto estado (p. ej. «Crítico»). \\
\hline
RN-031 & Origen de las sanciones & \textbf{SI} el sistema identifica sanciones, \textbf{ENTONCES} solo las toma de menciones explícitas en los documentos ingestados de la empresa. \textbf{SE APLICA A} el análisis de sanciones. \textbf{SE BASA EN} que el sistema no consulta fuentes externas. \textbf{Y TIENE LA SIGUIENTE EXCEPCIÓN} si la sanción no tiene monto, se conserva como «no cuantificada» y nunca se asume cero (RN-032). \\
\hline
RN-032 & Cálculo del puntaje ESG & \textbf{SI} se asignan estados a los códigos GRI de una empresa y un año, \textbf{ENTONCES} el sistema calcula el puntaje ESG con la fórmula OK = 100, Baja sustancia = 50 y Sub-reportado = 0. \textbf{SE APLICA A} el reporte de prospección. \textbf{SE BASA EN} los estados validados. \textbf{Y TIENE LA SIGUIENTE EXCEPCIÓN} si no se detectó ningún código, el puntaje se muestra como «no disponible», nunca como cero. \\
\hline
RN-033 & Fundamentación en el corpus & \textbf{SI} el Administrador consulta al asistente, \textbf{ENTONCES} el sistema genera la respuesta exclusivamente con base en los documentos ingestados. \textbf{SE APLICA A} las consultas en lenguaje natural. \textbf{SE BASA EN} el acuerdo con el PO (IA delimitada a la base del cliente). \textbf{Y TIENE LA SIGUIENTE EXCEPCIÓN} no se emplea conocimiento externo al corpus. \\
\hline
RN-034 & Trazabilidad de la información entregada & \textbf{SI} el asistente entrega información, \textbf{ENTONCES} toda afirmación debe ser atribuible a un documento del corpus, citando documento, empresa y año. \textbf{SE APLICA A} las respuestas del asistente. \textbf{SE BASA EN} la verificabilidad de la información. \textbf{Y TIENE LA SIGUIENTE EXCEPCIÓN} una afirmación sin respaldo trazable no es información válida del sistema. \\
\hline
RN-035 & Declaración explícita ante ausencia de información & \textbf{SI} el corpus no contiene información suficiente, \textbf{ENTONCES} el asistente lo declara explícitamente en lugar de inventar una respuesta. \textbf{SE APLICA A} las respuestas del asistente. \textbf{SE BASA EN} la honestidad de la respuesta. \textbf{Y TIENE LA SIGUIENTE EXCEPCIÓN} la ausencia de respaldo nunca se sustituye por datos generados o atribuidos a fuentes inexistentes. \\
\hline
RN-036 & Restricción de dominio del asistente & \textbf{SI} la consulta se refiere a un tema fuera de sostenibilidad empresarial, indicadores GRI, sanciones o el contenido ingestado, \textbf{ENTONCES} el asistente declina explícitamente. \textbf{SE APLICA A} todas las consultas. \textbf{SE BASA EN} el alcance funcional acordado. \textbf{Y TIENE LA SIGUIENTE EXCEPCIÓN} declina aunque el modelo tenga conocimiento general suficiente para responderla. \\
\hline
RN-037 & Requisitos y contexto de la consulta & \textbf{SI} el Administrador formula una consulta, \textbf{ENTONCES} debe haber seleccionado sector, empresa y año, y la respuesta se limita a ese contexto. \textbf{SE APLICA A} las consultas al asistente. \textbf{SE BASA EN} el contexto obligatorio de la consulta. \textbf{Y TIENE LA SIGUIENTE EXCEPCIÓN} sin esos datos el sistema no permite realizar la consulta. \\
\hline
RN-038 & Precondición del reporte de prospección & \textbf{SI} se genera un reporte de prospección, \textbf{ENTONCES} la empresa debe contar con al menos un documento indexado. \textbf{SE APLICA A} la generación de reportes. \textbf{SE BASA EN} la disponibilidad de resultados analizables. \textbf{Y TIENE LA SIGUIENTE EXCEPCIÓN} si no hay documentos, el sistema bloquea la generación e informa el motivo. \\
\hline
RN-039 & Contenido del reporte de prospección & \textbf{SI} se genera un reporte, \textbf{ENTONCES} consolida, para una empresa y un año, el estado de cada código GRI con su cita, las sanciones identificadas (con monto o «no cuantificada») y un resumen ejecutivo con el puntaje ESG. \textbf{SE APLICA A} la generación de reportes. \textbf{SE BASA EN} el resultado almacenado y validado. \textbf{Y TIENE LA SIGUIENTE EXCEPCIÓN} no se incluye información sin respaldo en el corpus. \\
\hline
RN-040 & Determinismo y generación sin IA & \textbf{SI} se genera un reporte, \textbf{ENTONCES} su contenido se construye exclusivamente con los resultados ya almacenados y validados, sin invocar al servicio de IA. \textbf{SE APLICA A} la generación de reportes. \textbf{SE BASA EN} la reproducibilidad del resultado. \textbf{Y TIENE LA SIGUIENTE EXCEPCIÓN} dos generaciones sobre los mismos datos producen el mismo resultado. \\
\hline
RN-041 & Inmutabilidad y visualización del reporte & \textbf{SI} un reporte es generado, \textbf{ENTONCES} este no puede modificarse ni eliminarse, y su historial es visible y descargable por cualquier Administrador. \textbf{SE APLICA A} los reportes de prospección. \textbf{SE BASA EN} la integridad del entregable. \textbf{Y TIENE LA SIGUIENTE EXCEPCIÓN} no se genera un segundo reporte de la misma empresa y año. \\
\hline
RN-042 & Eventos auditados & \textbf{SI} ocurre un evento sensible (inicio de sesión, cambio de estado o rol de cuenta, ingesta, rechazo de documento o generación de reporte), \textbf{ENTONCES} el sistema lo registra automáticamente. \textbf{SE APLICA A} toda la plataforma. \textbf{SE BASA EN} la trazabilidad de las operaciones. \textbf{Y TIENE LA SIGUIENTE EXCEPCIÓN} no se permite excluir ninguna de esas acciones del registro. \\
\hline
RN-043 & Contenido del registro de auditoría & \textbf{SI} el sistema registra un evento, \textbf{ENTONCES} consigna la cuenta que originó la acción, la fecha y hora (UTC), el tipo de acción y el resultado. \textbf{SE APLICA A} los eventos auditados. \textbf{SE BASA EN} la identificación del responsable. \textbf{Y TIENE LA SIGUIENTE EXCEPCIÓN} en eventos anónimos (p. ej. inicio de sesión fallido de un correo inexistente) la cuenta puede ser nula. \\
\hline
RN-044 & Inmutabilidad de la auditoría & \textbf{SI} se registra un evento de auditoría, \textbf{ENTONCES} este no puede modificarse ni eliminarse; el registro es de solo inserción. \textbf{SE APLICA A} la bitácora de auditoría. \textbf{SE BASA EN} la integridad del registro. \textbf{Y TIENE LA SIGUIENTE EXCEPCIÓN} la consulta es de solo lectura y el registro se conserva durante todo el periodo de operación. \\
\hline
RN-045 & Acceso a la auditoría & \textbf{SI} una cuenta distinta del SuperAdmin intenta acceder a la auditoría, \textbf{ENTONCES} el sistema deniega el acceso y registra el intento. \textbf{SE APLICA A} la bitácora de auditoría. \textbf{SE BASA EN} la restricción por rol. \textbf{Y TIENE LA SIGUIENTE EXCEPCIÓN} solo el SuperAdmin puede consultar y filtrar los eventos. \\
\hline
\end{longtable}
